\section{Formal Consequence Analysis}
\label{sec:formalization}
\label{sec:attack-path}

\paragraph{Decision input.}
\label{sec:implementation-mapping}
The bridge is a provenance-checked transformation component for this case
study. It checks the recorded ordered rules, initiating identity, and
historical matcher provenance, and combines independent finite
specification evaluation with the archived implementation observation:
\[
 (D_{\mathrm{spec}}(P,I),D_{\mathrm{impl}}^{\mathrm{obs}}(P,I))
       =(\mathsf{Deny},\mathsf{Allow}).
\]
The supported evaluator grammar covers the literal and wildcard patterns
needed for this case; it is not complete Python or \texttt{fnmatch}
semantics. The checked pair reconstructs only the scenario region of a
fixed consequence template. Protocol processes and queries remain
unchanged; symbolic role mapping remains an analysis choice.

\begin{figure}[htbp]
\centering
\begin{tabular}{c}
Implementation observation + specification evaluation \\
$\downarrow$ \\
Checked decision pair $\longrightarrow$ fixed symbolic scenario \\
$\downarrow$ \\
ProVerif analysis under explicit assumptions
\end{tabular}
\caption{Case-specific decision transfer; the arrows do not establish
implementation refinement.}
\label{fig:implementation-mapping}
\end{figure}

Table~\ref{tab:implementation-mapping} gives the scenario facts.
ProVerif consumes these fixed inputs; it neither executes the matcher nor
discovers its inconsistency.

\begin{table}[htbp]
\centering
\small
\begin{tabular}{ll}
\hline
Checked decision & Symbolic table fact \\
\hline
Specification: Deny & \texttt{ScenarioSpecDeny(BuggyPolicy,aid\_B)} \\
Observation: Allow & \texttt{ScenarioImplAllow(BuggyPolicy,aid\_B)} \\
\hline
\end{tabular}
\caption{Decision values transferred into the fixed scenario.}
\label{tab:implementation-mapping}
\end{table}

\paragraph{Model and assumptions.}
The initiating peer $I$ has a genuine registered identity and an existing
receiving-user policy outside its control. The scenario assumes the same
fixed policy at the Provider and receiver, effective cryptographic checks,
and no compromise of the Provider, CA, or honest private keys.
The prescribed initiating-peer process does not represent arbitrary
adversary-controlled application behavior.

The Provider releases contact material; \texttt{PeerA} represents receiver
$R$, which issues tokens and accepts messages; \texttt{PeerB} represents
initiator $I$. A CA supplies certificates. Replicated processes use symbolic
cryptography, and the private \texttt{tls\_chat} channel abstracts
authenticated, confidential transport. The \texttt{ProviderReleased}
table abstracts handoff to receiver processing; \texttt{ActiveToken}
associates the participants, token, and sender's access-control public key.

\texttt{SpecIntendedDeny}$(I,R)$ records scenario denial without blocking
execution. \texttt{ProviderRelease}$(R,I)$ records release,
\texttt{TokenIssue}$(R,I,t)$ records wire-token issuance, and
\texttt{ChatAccept}$(R,I,t,m)$ records application-message acceptance.
The latter is distinct from \texttt{Accept}, which denotes token receipt
and verification in the diagnostic model.

\paragraph{Queries and results.}
The archived ProVerif~2.05 execution of the bridge-derived scenario reports
the two results in Table~\ref{tab:claim-formalization}. A reconstructed
witness refutes \texttt{not event(ChatAccept(...))}; acceptance is
existentially reachable under the supplied divergence. The non-injective
correspondence
\[
 \mathsf{ChatAccept}(R,I,t,m)
 \Longrightarrow \mathsf{TokenIssue}(R,I,t)
\]
requires prior issuance for the same participants and token, not intended
policy permission or a unique session. Issuance occurs in the witness;
there is no standalone issuance-reachability query in this model.

\begin{table}[htbp]
\centering
\begin{tabular}{ll}
\hline
Query in the derived-model run & Result \\
\hline
\texttt{ChatAccept} reachability & Reachable \\
\texttt{ChatAccept} $\Rightarrow$ \texttt{TokenIssue} & True \\
\hline
\end{tabular}
\caption{Archived formal query outcomes for the fixed D/A scenario.}
\label{tab:claim-formalization}
\label{tab:evaluation-runs}
\end{table}

The run is archived under \texttt{proofs/evidence/}\allowbreak
\texttt{authz-bridge-turepass-}\allowbreak\texttt{20260922T111657792091Z/}.
Its manifest binds the observation, generated model, verifier, and outputs.
It is separate from the historical preset-divergence run and the consumer
experiment; agreement of outcomes does not identify a common execution.

The earlier Stage~1 diagnostic retains the queried authentication
correspondences and token secrecy, while
$\mathsf{Accept}(R,I,t)\Rightarrow\mathsf{PolicyAllow}(R,I)$ is false.
Because that process never emits \texttt{PolicyAllow}, this is a
verification-coverage diagnostic, not formal detection of the matcher
defect. Those guarantees do not automatically transfer to the consequence
model.

\paragraph{Independent gate controls.}
The existing controls in Table~\ref{tab:gate-evaluation} isolate modeled
blocking positions. Their labels denote Provider/receiver permissions,
not specification/implementation decisions. They use a fixed private
token, unlike the fresh-token divergence model, and are not verified
repairs or refinements of it.

\begin{table}[htbp]
\centering
\begin{tabular}{llccc}
\hline
Provider & Receiver & Release & Issue & Accept \\
\hline
Allow & Allow & R & R & R \\
Deny & Allow & U & U & U \\
Allow & Deny & R & U & U \\
\hline
\end{tabular}
\caption{Independent gate controls: reachable (R) and unreachable (U).}
\label{tab:gate-evaluation}
\end{table}

In addition to the prior-issuance correspondence, the controls establish
\[
 \mathsf{ChatAccept}(R,I,t,m)\Rightarrow\mathsf{ChatSend}(R,I,t,m).
\]
Both correspondences are vacuous when acceptance is unreachable.
Provider denial blocks release and downstream events; receiver denial
blocks issuance and acceptance while release remains reachable.
The evidence remains in \texttt{proofs/evidence/}\allowbreak
\texttt{authz-chat-gates-}\allowbreak\texttt{20260918T023621017067Z/}.
